\hypertarget{appendix-db3aef8f2fd6}{}
\section*{Cubic Hermite interpolation}
\section{Cubic Hermite interpolation}\label{cubic-hermite-interpolation}

\textbf{Curriculum depth:} Developed · \textbf{First seen:} Phase 10 · \textbf{Latest development:} Phase 10

\subsection{Definition and mechanism}\label{definition-and-mechanism}

A cubic Hermite segment is determined by endpoint values and endpoint tangent data. With normalized \(t\in[0,1]\), its four basis polynomials blend the two endpoint positions and two tangents. Tangents are derivatives with respect to the parameter, so physical-time derivatives must be scaled by segment duration: a velocity in units per second is not interchangeable with a tangent per normalized parameter. Hermite curves are useful when a camera path or animation must pass through known positions while exposing direct control over its entry and exit directions. The endpoint constraints determine one cubic polynomial; continuity between adjacent segments still requires compatible endpoint tangents.

\subsection{Prerequisites}\label{prerequisites}

\begin{itemize}
\tightlist
\item
  \hyperref[appendix-0c2a5f69851e]{Interpolation versus approximation}
\item
  \href{}{Spline bases and curve evaluation (03-Visual-Computing-I, appendix-f28eeec97a3d)}
\end{itemize}

\subsection{Used by}\label{used-by}

\begin{itemize}
\tightlist
\item
  No downstream concept node is added in this stage.
\end{itemize}

\subsection{Full teaching explanation}\label{full-teaching-explanation}

\begin{itemize}
\tightlist
\item
  \href{Continuous\%20Notes/Computing\%20Foundations\%20-\%20From\%20Bits\%20to\%20C++\%20and\%20Systems\%20APIs\#Reconstructing\%20values\%20and\%20curves\%20between\%20samples}{Reconstructing values and curves between samples}
\end{itemize}

\subsection{Source coverage}\label{source-coverage}

Supplied lesson parts: 981.

\subsection{Future keywords}\label{future-keywords}

endpoint tangents, C1 continuity, parameter scaling
